\documentclass[10pt, letterpaper]{article}

% Packages:
\usepackage[
    ignoreheadfoot, % set margins without considering header and footer
    top=2 cm, % seperation between body and page edge from the top
    bottom=2 cm, % seperation between body and page edge from the bottom
    left=2 cm, % seperation between body and page edge from the left
    right=2 cm, % seperation between body and page edge from the right
    footskip=1.0 cm, % seperation between body and footer
    % showframe % for debugging 
]{geometry} % for adjusting page geometry
\usepackage{titlesec} % for customizing section titles
\usepackage{tabularx} % for making tables with fixed width columns
\usepackage{array} % tabularx requires this
\usepackage[dvipsnames]{xcolor} % for coloring text
\definecolor{primaryColor}{RGB}{0, 0, 0} % define primary color
\usepackage{enumitem} % for customizing lists
\usepackage{fontawesome5} % for using icons
\usepackage{amsmath} % for math
\usepackage[
    pdftitle={John Doe's CV},
    pdfauthor={John Doe},
    pdfcreator={LaTeX with RenderCV},
    colorlinks=true,
    urlcolor=primaryColor
]{hyperref} % for links, metadata and bookmarks
\usepackage[pscoord]{eso-pic} % for floating text on the page
\usepackage{calc} % for calculating lengths
\usepackage{bookmark} % for bookmarks
\usepackage{lastpage} % for getting the total number of pages
\usepackage{changepage} % for one column entries (adjustwidth environment)
\usepackage{paracol} % for two and three column entries
\usepackage{ifthen} % for conditional statements
\usepackage{needspace} % for avoiding page brake right after the section title
\usepackage{iftex} % check if engine is pdflatex, xetex or luatex

% Ensure that generate pdf is machine readable/ATS parsable:
\ifPDFTeX
    \input{glyphtounicode}
    \pdfgentounicode=1
    \usepackage[T1]{fontenc}
    \usepackage[utf8]{inputenc}
    \usepackage{lmodern}
\fi

\usepackage{charter}

% Some settings:
\raggedright
\AtBeginEnvironment{adjustwidth}{\partopsep0pt} % remove space before adjustwidth environment
\pagestyle{empty} % no header or footer
\setcounter{secnumdepth}{0} % no section numbering
\setlength{\parindent}{0pt} % no indentation
\setlength{\topskip}{0pt} % no top skip
\setlength{\columnsep}{0.15cm} % set column seperation
\pagenumbering{gobble} % no page numbering

\titleformat{\section}{\needspace{4\baselineskip}\bfseries\large}{}{0pt}{}[\vspace{1pt}\titlerule]

\titlespacing{\section}{
    % left space:
    -1pt
}{
    % top space:
    0.3 cm
}{
    % bottom space:
    0.2 cm
} % section title spacing

\renewcommand\labelitemi{$\vcenter{\hbox{\small$\bullet$}}$} % custom bullet points
\newenvironment{highlights}{
    \begin{itemize}[
        topsep=0.10 cm,
        parsep=0.10 cm,
        partopsep=0pt,
        itemsep=0pt,
        leftmargin=0 cm + 10pt
    ]
}{
    \end{itemize}
} % new environment for highlights


\newenvironment{highlightsforbulletentries}{
    \begin{itemize}[
        topsep=0.10 cm,
        parsep=0.10 cm,
        partopsep=0pt,
        itemsep=0pt,
        leftmargin=10pt
    ]
}{
    \end{itemize}
} % new environment for highlights for bullet entries

\newenvironment{onecolentry}{
    \begin{adjustwidth}{
        0 cm + 0.00001 cm
    }{
        0 cm + 0.00001 cm
    }
}{
    \end{adjustwidth}
} % new environment for one column entries

\newenvironment{twocolentry}[2][]{
    \onecolentry
    \def\secondColumn{#2}
    \setcolumnwidth{\fill, 4.5 cm}
    \begin{paracol}{2}
}{
    \switchcolumn \raggedleft \secondColumn
    \end{paracol}
    \endonecolentry
} % new environment for two column entries

\newenvironment{threecolentry}[3][]{
    \onecolentry
    \def\thirdColumn{#3}
    \setcolumnwidth{, \fill, 4.5 cm}
    \begin{paracol}{3}
    {\raggedright #2} \switchcolumn
}{
    \switchcolumn \raggedleft \thirdColumn
    \end{paracol}
    \endonecolentry
} % new environment for three column entries

\newenvironment{header}{
    \setlength{\topsep}{0pt}\par\kern\topsep\centering\linespread{1.5}
}{
    \par\kern\topsep
} % new environment for the header

\newcommand{\placelastupdatedtext}{% \placetextbox{<horizontal pos>}{<vertical pos>}{<stuff>}
  \AddToShipoutPictureFG*{% Add <stuff> to current page foreground
    \put(
        \LenToUnit{\paperwidth-2 cm-0 cm+0.05cm},
        \LenToUnit{\paperheight-1.0 cm}
    ){\vtop{{\null}\makebox[0pt][c]{
             \small\color{gray}\textit{Last updated in October 2025}\hspace{\widthof{Last updated in October 2025}}
    }}}%
  }%
}%

% save the original href command in a new command:
\let\hrefWithoutArrow\href

% new command for external links:


\begin{document}
    %% New CV content starts here
    \newcommand{\AND}{\unskip
        \cleaders\copy\ANDbox\hskip\wd\ANDbox
        \ignorespaces
    }
    \newsavebox\ANDbox
    \sbox\ANDbox{$|$}

    \placelastupdatedtext

    % Header
    \begin{header}
        \textbf{\fontsize{24 pt}{24 pt}\selectfont Aaron Feng}

        \vspace{5 pt}

        \normalsize
        \mbox{{\color{black}\footnotesize\faMapMarker*}\hspace*{0.13cm}Del Mar, CA}%
        \kern 5.0 pt%
        \AND%
        \kern 5.0 pt%
        \mbox{\hrefWithoutArrow{mailto:zhf004@ucsd.edu}{\color{black}{\footnotesize\faEnvelope[regular]}\hspace*{0.13cm}zhf004@ucsd.edu}}%
        \kern 5.0 pt%
        \AND%
        \kern 5.0 pt%
        \mbox{\hrefWithoutArrow{tel:+13363377139}{\color{black}{\footnotesize\faPhone*}\hspace*{0.13cm}(336) 337\,7139}}%
        \kern 5.0 pt%
        \AND%
        \kern 5.0 pt%
        \mbox{\hrefWithoutArrow{https://github.com/aaronzhfeng}{\color{black}{\footnotesize\faGithub}\hspace*{0.13cm}aaronzhfeng}}%
        \kern 5.0 pt%
        \AND%
        \kern 5.0 pt%
        \mbox{\hrefWithoutArrow{https://zhf004.github.io}{\color{black}{\footnotesize\faGlobe}\hspace*{0.13cm}zhf004.github.io}}%
    \end{header}

    \vspace{5 pt - 0.3 cm}

    % Education
    \section{Education}

    \begin{twocolentry}{\textit{Sept.\@ 2022 -- Jun.\@ 2026 (expected)}}
        \textbf{University of California, San Diego (UCSD)}

        \textit{B.S. Data Science \& B.S. Probability \& Statistics}
    \end{twocolentry}

    \vspace{0.10 cm}
    \begin{onecolentry}
        \begin{highlights}
            \item GPA: 3.83 (Major GPAs: 3.94 in Data Science, 3.93 in Statistics)
            \item Honors: Provost Honors (multiple quarters)【250326654186019†L220-L224】
            \item Graduate-level Coursework: DSC\,291 (Topics in Data Science), MATH\,287C (Advanced Time Series Analysis), DSC\,240 (Machine Learning), MATH\,282B (Applied Statistics\,II), MATH\,287A (Time Series Analysis), MATH\,278C (Seminar in Optimization), MATH\,282A (Applied Statistics\,I)
            \item Relevant Coursework: Machine Learning, Deep Learning, Natural Language Processing, Mathematical Statistics, Stochastic Processes, Data Management, Linear Algebra, Analysis
        \end{highlights}
    \end{onecolentry}

    % Publication
    \section{Publication}

    \begin{onecolentry}
        \textbf{ArcMemo: Abstract Reasoning Composition with Lifelong LLM Memory} \\
        \textit{Matthew Ho, Chen Si, Zhaoxiang Feng, et al.} (arXiv preprint, Oct.\@\,2025) \\
        \href{https://arxiv.org/abs/2509.04439}{Paper} \\
        Our work introduces \emph{ArcMemo}, a concept-level memory framework for large language models. By distilling reusable abstractions from reasoning traces, ArcMemo enables continual learning at test time without weight updates and yields a 7.5\% relative gain over a strong no-memory baseline on the ARC-AGI benchmark【544438269244232†L20-L34】. The memory design outperforms instance-level approaches across inference compute scales【544438269244232†L135-L139】.
    \end{onecolentry}

    % Research Experience
    \section{Research Experience}

    % Q-lab Research Assistant
    \begin{twocolentry}{\textit{Jan.\@ 2025 – Present}}
        \textbf{Q-lab} \\
        \textit{Research Assistant}
    \end{twocolentry}

    \vspace{0.10 cm}
    \begin{onecolentry}
        \begin{highlights}
            \item Led the development of \emph{ArcMemo}, a lifelong memory module for LLM reasoning. Built a pipeline to extract and abstract concepts from solution traces and integrate them into prompts, enabling continual improvement across tasks【544438269244232†L20-L34】.
            \item Evaluated ArcMemo on the ARC‑AGI benchmark; achieved a 7.5\% relative improvement over a strong no‑memory baseline and demonstrated performance scaling with inference compute【544438269244232†L135-L139】.
            \item Designed dynamic feature detection routines and helper functions to generate transformation hints for compositional puzzles, and optimized asynchronous API workflows for efficient model evaluation.
        \end{highlights}
    \end{onecolentry}

    \vspace{0.20 cm}

    % USIPC Research Assistant
    \begin{twocolentry}{\textit{Sept.\@ 2023 – Jun.\@ 2024}}
        \textbf{U.S. Immigration Policy Center} \\
        \textit{Research Assistant}
    \end{twocolentry}

    \vspace{0.10 cm}
    \begin{onecolentry}
        \begin{highlights}
            \item Utilized Python and R to analyze historical naturalization data and forecast long‑term immigration trends, informing policy briefs for the 2024 election cycle.
            \item Developed reproducible data pipelines and visualization dashboards to communicate findings to non‑technical stakeholders.
        \end{highlights}
    \end{onecolentry}

    \vspace{0.20 cm}

    % Wang Lab Student Researcher (placeholder)
    \begin{twocolentry}{\textit{Jun.\@ 2025 – Present}}
        \textbf{Wang Lab} \\
        \textit{Student Researcher}
    \end{twocolentry}

    \vspace{0.10 cm}
    \begin{onecolentry}
        \begin{highlights}
            \item Working on computational chemistry projects exploring graph neural networks and transformer models; details forthcoming.
        \end{highlights}
    \end{onecolentry}

    \vspace{0.20 cm}

    % HDSI Student Researcher – Hao Zhang
    \begin{twocolentry}{\textit{Sept.\@ 2025 – Present}}
        \textbf{Halıcıo\u{g}lu Data Science Institute, UC San Diego} \\
        \textit{Student Researcher} — Advisors: Prof.~Hao Zhang
    \end{twocolentry}

    \vspace{0.10 cm}
    \begin{onecolentry}
        \begin{highlights}
            \item Senior Capstone: “Open LLM Training, Inference, and Infrastructure” (ongoing). Building scalable pipelines for training open‑source large language models, benchmarking inference frameworks, and exploring alignment strategies.
        \end{highlights}
    \end{onecolentry}

    \vspace{0.20 cm}

    % HDSI Student Researcher – Tian Ma
    \begin{twocolentry}{\textit{Sept.\@ 2025 – Present}}
        \textbf{Halıcıo\u{g}lu Data Science Institute, UC San Diego} \\
        \textit{Student Researcher} — Advisors: Prof.~Tian Ma
    \end{twocolentry}

    \vspace{0.10 cm}
    \begin{onecolentry}
        \begin{highlights}
            \item Senior Capstone: “Quantifying the credibility of large language model outputs” (ongoing). Investigating metrics and algorithms to measure factual consistency and reliability in LLM responses.
        \end{highlights}
    \end{onecolentry}

    % Teaching Experience
    \section{Teaching Experience}

    \begin{twocolentry}{\textit{Fall 2024}}
        \textbf{Reader — MATH 180C: Introduction to Stochastic Processes II} \\
        University of California, San Diego
    \end{twocolentry}

    \vspace{0.10 cm}
    \begin{onecolentry}
        \begin{highlights}
            \item Graded problem sets and exams, ensuring detailed feedback on Markov chains and stochastic calculus topics.
            \item Held weekly office hours to clarify course material and support student comprehension.
        \end{highlights}
    \end{onecolentry}

    \vspace{0.15 cm}

    \begin{twocolentry}{\textit{Winter 2025}}
        \textbf{Reader — MATH 185: Introduction to Computational Statistics} \\
        University of California, San Diego
    \end{twocolentry}

    \vspace{0.10 cm}
    \begin{onecolentry}
        \begin{highlights}
            \item Evaluated programming assignments and exams covering statistical simulation, Monte Carlo methods, and estimation techniques.
            \item Collaborated with instructors to provide clarifications on computational projects.
        \end{highlights}
    \end{onecolentry}

    \vspace{0.15 cm}

    \begin{twocolentry}{\textit{Spring 2025}}
        \textbf{Reader — MATH 180C: Introduction to Stochastic Processes II} \\
        University of California, San Diego
    \end{twocolentry}

    \vspace{0.10 cm}
    \begin{onecolentry}
        \begin{highlights}
            \item Provided detailed feedback on assignments focusing on Poisson processes and discrete‑time martingales.
            \item Assisted with exam proctoring and maintained grade records.
        \end{highlights}
    \end{onecolentry}

    \vspace{0.15 cm}

    \begin{twocolentry}{\textit{Spring 2025}}
        \textbf{Reader — MATH 181A: Introduction to Mathematical Statistics I} \\
        University of California, San Diego
    \end{twocolentry}

    \vspace{0.10 cm}
    \begin{onecolentry}
        \begin{highlights}
            \item Graded problem sets and exams on estimation, hypothesis testing, and confidence intervals.
            \item Supported the instructor by addressing student questions via discussion forums.
        \end{highlights}
    \end{onecolentry}

    \vspace{0.15 cm}

    \begin{twocolentry}{\textit{Summer Session I 2025}}
        \textbf{Tutor — DSC 40A: Theoretical Foundations of Data Science I} \\
        University of California, San Diego
    \end{twocolentry}

    \vspace{0.10 cm}
    \begin{onecolentry}
        \begin{highlights}
            \item Conducted tutoring sessions covering set theory, probability, and algorithmic thinking for early‑career data science students.
            \item Led review sessions and prepared supplementary materials to aid student learning.
        \end{highlights}
    \end{onecolentry}

    \vspace{0.15 cm}

    \begin{twocolentry}{\textit{Summer Session II 2025}}
        \textbf{Reader — MATH 181A: Introduction to Mathematical Statistics I} \\
        University of California, San Diego
    \end{twocolentry}

    \vspace{0.10 cm}
    \begin{onecolentry}
        \begin{highlights}
            \item Assisted with grading and feedback for advanced statistical inference coursework.
        \end{highlights}
    \end{onecolentry}

    \vspace{0.15 cm}

    \begin{twocolentry}{\textit{Fall 2025}}
        \textbf{Reader — MATH 180A: Introduction to Probability (Sections B00 \& C00)} \\
        University of California, San Diego
    \end{twocolentry}

    \vspace{0.10 cm}
    \begin{onecolentry}
        \begin{highlights}
            \item Managed grading for two lecture sections, covering probability spaces, random variables, and limit theorems.
            \item Provided consultation during office hours on problem‑solving strategies.
        \end{highlights}
    \end{onecolentry}

    % Selected Projects
    \section{Selected Projects}

    % 3D-Enhanced Graph Reaction Prediction
    \begin{twocolentry}{\textit{\href{https://aaronzhfeng.github.io/projects/3d-enhanced-graph-reaction}{3D‑Enhanced Graph Reaction Prediction}}}
        \textbf{3D‑Enhanced Graph Reaction Prediction}
    \end{twocolentry}

    \vspace{0.10 cm}
    \begin{onecolentry}
        \begin{highlights}
            \item Designed a graph neural network for reaction outcome prediction combining 2D molecular graph encoders with 3D geometry‑aware message passing via a contrastive InfoNCE objective【423483429983654†L17-L28】.
            \item Pre‑trained 2D and 3D encoders on the QM9 dataset to infuse spatial knowledge and fine‑tuned on the Open Reaction Database (ORD); observed that incorporating 3D features improves representations for stereochemically constrained reactions【423483429983654†L17-L28】.
            \item Integrated the graph encoder with an LSTM decoder to generate product SMILES, demonstrating improved accuracy on reactions with significant geometric constraints.
        \end{highlights}
    \end{onecolentry}

    \vspace{0.2 cm}

    % ChemTransformer
    \begin{twocolentry}{\textit{\href{https://aaronzhfeng.github.io/projects/chemtransformer}{ChemTransformer}}}
        \textbf{ChemTransformer}
    \end{twocolentry}

    \vspace{0.10 cm}
    \begin{onecolentry}
        \begin{highlights}
            \item Built a transformer‑based sequence‑to‑sequence model for chemical reaction prediction, leveraging SMILES and SELFIES tokenization and novel compound‑based relative positional encodings【831730747525686†L18-L34】.
            \item Explored multiple model configurations on the Open Reaction Dataset (ORD) and found that traditional token‑level positional encoding outperformed compound‑level variants, while SELFIES improved output validity【831730747525686†L30-L38】.
            \item Achieved competitive sequence‑level accuracy despite training challenges, highlighting the potential of transformer architectures for large‑scale reaction prediction【831730747525686†L30-L40】.
        \end{highlights}
    \end{onecolentry}

    % Honors & Awards
    \section{Honors \& Awards}

    \begin{onecolentry}
        \begin{highlights}
            \item Provost Honors (Fall 2022 – Spring 2025)【250326654186019†L220-L224】.
        \end{highlights}
    \end{onecolentry}

    % Technical Skills
    \section{Technical Skills}

    \begin{onecolentry}
        \textbf{Languages/Frameworks:} Python, Java, R, SQL, JavaScript, HTML \\
        \textbf{ML \& Data Tools:} PyTorch, scikit‑learn, Pandas, Dask, Spark, AWS \\
        \textbf{NLP/LLM:} Transformers, LLM fine‑tuning, Retrieval‑Augmented Generation, agent building, GraphRAG, OpenAI/Anthropic APIs (sync/async requests) \\
        \textbf{Systems/Tooling:} Distributed training, Weights \& Biases, Docker, Linux, SSH
    \end{onecolentry}

\end{document}
        \end{onecolentry}

        \vspace{0.2 cm}

        \begin{onecolentry}
            The boilerplate content was inspired by \href{https://github.com/dnl-blkv/mcdowell-cv}{Gayle McDowell}.
        \end{onecolentry}


    
    \section{Quick Guide}

    \begin{onecolentry}
        \begin{highlightsforbulletentries}


        \item Each section title is arbitrary and each section contains a list of entries.

        \item There are 7 unique entry types: \textit{BulletEntry}, \textit{TextEntry}, \textit{EducationEntry}, \textit{ExperienceEntry}, \textit{NormalEntry}, \textit{PublicationEntry}, and \textit{OneLineEntry}.

        \item Select a section title, pick an entry type, and start writing your section!

        \item \href{https://docs.rendercv.com/user_guide/}{Here}, you can find a comprehensive user guide for RenderCV.


        \end{highlightsforbulletentries}
    \end{onecolentry}

    \section{Education}



        
        \begin{twocolentry}{
            Sept 2000 – May 2005
        }
            \textbf{University of Pennsylvania}, BS in Computer Science\end{twocolentry}

        \vspace{0.10 cm}
        \begin{onecolentry}
            \begin{highlights}
                \item GPA: 3.9/4.0 (\href{https://example.com}{a link to somewhere})
                \item \textbf{Coursework:} Computer Architecture, Comparison of Learning Algorithms, Computational Theory
            \end{highlights}
        \end{onecolentry}



    
    \section{Experience}



        
        \begin{twocolentry}{
            June 2005 – Aug 2007
        }
            \textbf{Software Engineer}, Apple -- Cupertino, CA\end{twocolentry}

        \vspace{0.10 cm}
        \begin{onecolentry}
            \begin{highlights}
                \item Reduced time to render user buddy lists by 75\% by implementing a prediction algorithm
                \item Integrated iChat with Spotlight Search by creating a tool to extract metadata from saved chat transcripts and provide metadata to a system-wide search database
                \item Redesigned chat file format and implemented backward compatibility for search
            \end{highlights}
        \end{onecolentry}


        \vspace{0.2 cm}

        \begin{twocolentry}{
            June 2003 – Aug 2003
        }
            \textbf{Software Engineer Intern}, Microsoft -- Redmond, WA\end{twocolentry}

        \vspace{0.10 cm}
        \begin{onecolentry}
            \begin{highlights}
                \item Designed a UI for the VS open file switcher (Ctrl-Tab) and extended it to tool windows
                \item Created a service to provide gradient across VS and VS add-ins, optimizing its performance via caching
                \item Built an app to compute the similarity of all methods in a codebase, reducing the time from $\mathcal{O}(n^2)$ to $\mathcal{O}(n \log n)$
                \item Created a test case generation tool that creates random XML docs from XML Schema
                \item Automated the extraction and processing of large datasets from legacy systems using SQL and Perl scripts
            \end{highlights}
        \end{onecolentry}



    
    \section{Publications}



        
        \begin{samepage}
            \begin{twocolentry}{
                Jan 2004
            }
                \textbf{3D Finite Element Analysis of No-Insulation Coils}
            \end{twocolentry}

            \vspace{0.10 cm}
            
            \begin{onecolentry}
                \mbox{Frodo Baggins}, \mbox{\textbf{\textit{John Doe}}}, \mbox{Samwise Gamgee}

                \vspace{0.10 cm}
                
        \href{https://doi.org/10.1109/TASC.2023.3340648}{10.1109/TASC.2023.3340648}
        \end{onecolentry}
        \end{samepage}


    
    \section{Projects}



        
        \begin{twocolentry}{
            \href{https://github.com/sinaatalay/rendercv}{github.com/name/repo}
        }
            \textbf{Multi-User Drawing Tool}\end{twocolentry}

        \vspace{0.10 cm}
        \begin{onecolentry}
            \begin{highlights}
                \item Developed an electronic classroom where multiple users can simultaneously view and draw on a "chalkboard" with each person's edits synchronized
                \item Tools Used: C++, MFC
            \end{highlights}
        \end{onecolentry}


        \vspace{0.2 cm}

        \begin{twocolentry}{
            \href{https://github.com/sinaatalay/rendercv}{github.com/name/repo}
        }
            \textbf{Synchronized Desktop Calendar}\end{twocolentry}

        \vspace{0.10 cm}
        \begin{onecolentry}
            \begin{highlights}
                \item Developed a desktop calendar with globally shared and synchronized calendars, allowing users to schedule meetings with other users
                \item Tools Used: C\#, .NET, SQL, XML
            \end{highlights}
        \end{onecolentry}


        \vspace{0.2 cm}

        \begin{twocolentry}{
            2002
        }
            \textbf{Custom Operating System}\end{twocolentry}

        \vspace{0.10 cm}
        \begin{onecolentry}
            \begin{highlights}
                \item Built a UNIX-style OS with a scheduler, file system, text editor, and calculator
                \item Tools Used: C
            \end{highlights}
        \end{onecolentry}



    
    \section{Technologies}



        
        \begin{onecolentry}
            \textbf{Languages:} C++, C, Java, Objective-C, C\#, SQL, JavaScript
        \end{onecolentry}

        \vspace{0.2 cm}

        \begin{onecolentry}
            \textbf{Technologies:} .NET, Microsoft SQL Server, XCode, Interface Builder
        \end{onecolentry}


    

\end{document}